\subsection{Simplicial complexes}

Recall that a {\emphcolor simplex} in $\bR^N$ is a subspace,
$$ [v_0,\dots,v_n] = \set{\sum_i\lambda_iv_i}[\sum_i\lambda_i=1,\ 0\leq\lambda_i\in\bR] , $$
where $v_i\in\bR^N$ are {\emphcolor geometrically independent} (if $\sum_i\lambda_iv_i=0$ and $\sum_i\lambda_i=0$ then $\lambda_i=0$).
The {\emphcolor faces} of a simplex $[v_0,\dots,v_n]$ are the simplices generated by a subset of $\set{v_0,\dots,v_n}$.
The {\emphcolor standard $n$-simplex} is the one generated by the standard basis vectors of $\bR^{n+1}$:
$$ \abs{\Delta^n} = [e_0,\dots,e_n] = \set{(x_0,\dots,x_n)\in\bR^{n+1}}[\sum_ix_i=1,\ x_i\geq0] \subseteq \bR^{n+1} . $$

A {\emphcolor simplicial complex} in $\bR^N$ is a collection of simplices $K$ such that the face of any simplex in $K$ is also in $K$, and the intersection of any two
simplices in $K$ is a face of both of them.
Given a simplical complex $K$, we can define its {\emphcolor geometric realization} $\abs K$ to be the union of all the simplices in $K$.
A subset $F\subseteq\abs K$ is closed iff $F\cap\sigma$ is closed in the simplex $\sigma$ for all $\sigma\in K$.

Simplicial complexes are a good starting point, but they are insufficient for our purposes.
A straightforward first generalization is as follows.

\bdefn[title={Abstract simplicial complexes}]

    An {\emphcolor abstract simplicial complex} is a collection of sets $K=\set{X^n}_{n<\omega}$, where $X^k$ is a collection of subsets of $X^0$
    of cardinality $k+1$.
    Furthermore, it must satisfy the condition that any subset of $X^k$ of size $j+1$ belongs to $X^j$.

    Elements of $X^0$ are called the {\emphcolor vertices} of the complex.

\edefn

Alternatively an abstract simplicial complex is a set $K$ of finite subsets of $X^0$ which is closed under taking subsets.
Then $X^n$ is just the collection of sets in $K$ of cardinality $n+1$.

In a simplicial complex of $\bR^N$ we required also that the intersection of simplices be a simplex, but this is because we were dealing with simplices
embedded in an ambient space.

The standard $n$-simplex $\abs{\Delta^n}$ gives rise to the abstract simplicial complex of non-empty subsets of $\set{0,\dots,n}$.

\bdefn[title={Face maps}]

    We will assume that an abstract simplicial complex $K$ comes equipped with a total ordering on its vertices.
    Then for an $n$-simplex $\set{v_0,\dots,v_n}\in K$ we order its vertices and write $[v_0,\dots,v_n]$ to mean $v_0<\cdots<v_n$.

    We further define the {\emphcolor face maps} of the simplicial complex $K$ to be maps $d_0,\dots,d_n\colon X^n\to X^{n-1}$ given by
    $$ d_i[v_0,\dots,v_n] = [v_0,\dots,\hat v_i,\dots,v_n] $$
    where the hat means the vertex was removed.

\edefn

For each $n$ we have maps $d_i$ for $i\leq n$, to differentiate the face maps between simplices of differing dimension we could write $d_i^n$ instead of merely $d_i$,
but this is an overload of information that is not necessary.

There is a fundamental relation between face maps which is immediately clear:
$$ d_id_j = d_{j-1}d_i,\qquad \hbox{for $i<j$} . \eqcnt[faceid] $$
Indeed, both sides map $[v_0,\dots,v_n]$ to $[v_0,\dots,\hat v_i,\dots,\hat v_j,\dots,v_n]$.

\bdefn[title={Simplicial maps}]

    If $K,L$ are simplicial complexes, a {\emphcolor simplicial map} $f\colon K\to L$ is a function from vertices of $K$ to vertices of $L$ such that if
    $\set{v_0,\dots,v_n}\in K$ then $\set{f(v_i)}_i\in L$.
    Note that a simplicial map need not be injective.

\edefn

